\section{Model and method}
\label{sec:model}

\subsection{The classical kicked rotor and the standard map}

We use scaled units in which the kick period and the moment of inertia are both 1, so the Hamiltonian is
\cref{eq:hamiltonian}. Hamilton's equations are $\dot\theta = p$ and
$\dot p = K\sin\theta\,\sum_n\delta(t-n)$. Between kicks the momentum is constant and the angle grows
linearly. Across a kick the angle is continuous and the momentum jumps by $K\sin\theta$. Writing
$(\theta_n, p_n)$ for the state just before kick $n$, one period is
\begin{equation}
  p_{n+1} = p_n + K\sin\theta_n, \qquad \theta_{n+1} = \theta_n + p_{n+1} \pmod{2\pi}.
  \label{eq:standard_map}
\end{equation}
This is the Chirikov standard map \cite{Chirikov1979}. Its Jacobian determinant is 1, so it preserves
phase-space area, as any Hamiltonian flow must. Solving for the earlier state gives the inverse map,
$\theta_n = \theta_{n+1} - p_{n+1}$ and $p_n = p_{n+1} - K\sin\theta_n$, which we use for the time-reversal
test.

The map is unchanged by $p \to p + 2\pi$, so phase space can be drawn on a torus,
$\theta, p \in [0, 2\pi)$, and we do that for the phase portraits. In the dynamics, however, the momentum is
never wrapped, because its unbounded growth is exactly what we want to measure. Small displacements obey the
linearized map $\delta p' = \delta p + K\cos\theta\,\delta\theta$, $\delta\theta' = \delta\theta + \delta p'$.
Their average logarithmic growth rate per kick is the Lyapunov exponent $\lambda$, which we compute by
renormalizing the displacement after every kick.

\subsection{Quantization and the Floquet operator}

Because $\theta$ is an angle, the wave function is periodic, and $\hat p = -i\hbar\,\partial_\theta$ has
eigenvalues $p = \hbar m$ with integer $m$ and eigenfunctions $\ee^{im\theta}/\sqrt{2\pi}$. Here $\hbar$ is the
\emph{effective} Planck constant: in scaled units it is the ratio of Planck's constant to the rotor's action
scale, so it measures how quantum the motion is, and $\hbar\to 0$ is the classical limit.

Integrating the Schr\"odinger equation $i\hbar\,\partial_t\psi = H\psi$ across a kick, the
$\delta$-function term dominates and the kinetic term contributes nothing, so the kick multiplies the wave
function by a phase, $\psi(\theta) \to \ee^{-iK\cos\theta/\hbar}\,\psi(\theta)$, exactly. Between kicks the
rotor is free and each momentum component picks up the phase $\ee^{-i\hbar m^2/2}$. The evolution from just
before one kick to just before the next is therefore given by the one-period (Floquet) operator
\begin{equation}
  \hat U = \ee^{-i\hat p^2/2\hbar}\,\ee^{-iK\cos\hat\theta/\hbar}.
  \label{eq:floquet}
\end{equation}
Its matrix elements in the momentum basis follow from the Jacobi-Anger expansion of the kick,
\begin{equation}
  \ee^{-ik\cos\theta} = \sum_r (-i)^r J_r(k)\,\ee^{ir\theta}
  \quad\Longrightarrow\quad
  \langle m'|\hat U|m\rangle = \ee^{-i\hbar m'^2/2}\,(-i)^{m'-m}\,J_{m'-m}(K/\hbar).
  \label{eq:matrix_elements}
\end{equation}
Because $J_r(k)$ is tiny for $|r| > k$, one kick couples each momentum state to about $k = K/\hbar$
neighbors on either side. So the quantum dynamics depends on two numbers: $K$, which sets the classical
chaos, and $\hbar$, which sets the free phase and, through $k = K/\hbar$, the number of momentum states one
kick connects.

Expectation values in momentum space are sums, for example $\avg{p^2} = \hbar^2\sum_m m^2|\psi_m|^2$. To
compare the two theories we start them in matching states: the quantum rotor in the momentum eigenstate
$m = 0$, which is spread uniformly in angle with zero momentum, and the classical ensemble with $p_0 = 0$ and
$\theta_0$ uniform on $[0, 2\pi)$. After one kick both give exactly $\avg{p^2} = K^2/2$.

\subsection{Numerical method}

We store $\psi$ at $N$ equally spaced angles $\theta_j = 2\pi j/N$. A discrete Fourier transform, computed
with the fast Fourier transform (FFT) and normalized to be unitary, gives the momentum amplitudes $\psi_m$ for
$m = -N/2, \dots, N/2 - 1$. One period is then four steps:
multiply $\psi(\theta_j)$ by $\ee^{-iK\cos\theta_j/\hbar}$, transform to momentum space, multiply $\psi_m$ by
$\ee^{-i\hbar m^2/2}$, and transform back. The inverse step undoes these in reverse order with the conjugate
phases.

This looks like the split-step Fourier method, but it is exact here, with no time-step error. The kick is
diagonal in $\theta$, the free evolution is diagonal in $m$, and the two act at different times rather than
simultaneously, so there is nothing to split. The only approximations are the finite momentum range and
roundoff. The FFT treats momentum as periodic, so probability that reached $|m| = N/2$ would reappear on
the other side. We therefore track the \emph{edge probability}
\begin{equation}
  P_{\rm edge} = \sum_{|m| > 3N/8} |\psi_m|^2
  \label{eq:edge}
\end{equation}
throughout every run and require its largest value to stay below $10^{-10}$. Unless stated otherwise we use
$N = 2^{14} = 16384$ states, chosen by validation check 4 below. One kick then takes about a millisecond on a
laptop.

For the classical ensemble matched to the quantum state we place $M$ angles on the midpoint grid
$\theta_j = 2\pi(j + \tfrac12)/M$ with $p = 0$. Averaging over this grid is the trapezoid rule, which is
extremely accurate for smooth periodic functions, so early-time averages converge much faster than with random
sampling. To measure the diffusion rate $D(K)$ we also use a \emph{torus ensemble}, with $\theta$ and $p$ both
uniform at random (fixed seed). The uniform measure is invariant under the map, so this ensemble has no
start-up transient.

\subsection{Validation}
\label{sec:validation}

Before any physics run, the code had to pass five checks, written as 39 automated tests. Every tolerance in
them has a stated reason: machine precision, the statistical error of an ensemble, or a known physical
correction. \Cref{tab:validation} summarizes them, and \cref{app:validation} gives the details. As a test of
the tests, we planted six deliberate bugs (a wrong free phase, a wrong ordering of the momentum grid, a wrong
classical update, a wrong inverse map, a broken edge probability, and a kick strength off by 0.1\%). Each made
at least one test fail.

\begin{table}[!htbp]
\centering
\small
\begin{tabular}{@{}>{\raggedright}p{0.17\linewidth}>{\raggedright}p{0.36\linewidth}>{\raggedright\arraybackslash}p{0.39\linewidth}@{}}
\toprule
Check & What is tested & Result \\
\midrule
1. Norm & $\avg{\psi|\psi} = 1$ over 1000 kicks, forward and backward & drifts by about $3\times10^{-16}$ per kick
  ($2.9\times10^{-13}$ after 1000), within the FFT roundoff allowance $n\log_2 N\,\varepsilon$ \\
2. Resonance & $\hbar = 4\pi$: $\avg{p^2} = K^2n^2/2$ and $|\psi_m|^2 = J_m^2(nK/\hbar)$;
  $\hbar = 2\pi$: exact return every two kicks & largest relative error $2.9\times10^{-13}$ over 1000 kicks;
  distribution within $1.7\times10^{-15}$; return within $2.8\times10^{-15}$ \\
3. Early times & exact $\avg{p^2}$ after one and two kicks, quantum and classical; at three kicks the quantum
  value approaches the classical one like $\hbar^2$ & quantum within $10^{-14}$ of the formulas; classical
  within the ensemble's error; observed order in $\hbar$: 1.97 \\
4. Grid & $N = 2^{14}$ against $2^{15}$ at $\hbar = 0.25$, 1000 kicks; edge probability below $10^{-10}$ &
  edge probability $2.6\times10^{-27}$; doubling $N$ changes $\avg{p^2}$ only at roundoff level; $N = 2^{12}$
  fails (edge $3.1\times10^{-8}$) \\
5. Time reversal & $n$ kicks forward, then $n$ back & quantum returns within $2.9\times10^{-13}$ after 1000
  kicks; the classical return fails after about 38 kicks, as predicted by $\ln(1/\varepsilon)/\lambda$ \\
\bottomrule
\end{tabular}
\caption{The five validation checks ($\varepsilon = 2.2\times10^{-16}$ is the machine precision; $K = 5$
throughout). Details and the reason for every tolerance are in \cref{app:validation}.}
\label{tab:validation}
\end{table}
